% =============================================================================
% CESS Graduation Project Report
% EEG-Based Emotion Recognition Using Classical Machine Learning on the DEAP Dataset
% =============================================================================
\documentclass[12pt,a4paper]{report}

% --- Packages ---
\usepackage[utf8]{inputenc}
\usepackage[T1]{fontenc}
\usepackage{mathptmx}                    % Times New Roman
\usepackage[left=3cm,right=2.5cm,top=2.5cm,bottom=2.5cm]{geometry}
\usepackage{setspace}\onehalfspacing
\usepackage{graphicx}
\usepackage{booktabs}
\usepackage{longtable}
\usepackage{multirow}
\usepackage{array}
\usepackage{caption}
\usepackage{subcaption}
\usepackage{amsmath,amssymb}
\usepackage{algorithm}
\usepackage{algpseudocode}
\usepackage{hyperref}
\usepackage{cite}
\usepackage{enumitem}
\usepackage{xcolor}
\usepackage{listings}
\usepackage{float}
\usepackage{fancyhdr}

\hypersetup{
    colorlinks=true,
    linkcolor=black,
    citecolor=blue,
    urlcolor=blue,
}

\captionsetup[figure]{font=small,labelfont=bf,position=bottom}
\captionsetup[table]{font=small,labelfont=bf,position=top}

% --- Header/Footer ---
\pagestyle{fancy}
\fancyhf{}
\fancyhead[R]{\thepage}
\renewcommand{\headrulewidth}{0.4pt}

% =============================================================================
\begin{document}

% ---- TITLE PAGE ----
\begin{titlepage}
\centering
\vspace*{1cm}

{\large\bfseries Ain Shams University\\Faculty of Engineering}\\[0.3cm]
{\normalsize Computer Engineering and Software Systems (CESS) Program}\\[0.3cm]
{\small 1 Elsarayat St., Abbaseya, 11517 Cairo, Egypt}\\[2cm]

{\LARGE\bfseries Graduation Project Report}\\[1.5cm]

{\huge\bfseries EEG-Based Emotion Recognition\\[0.3cm]
Using Classical Machine Learning\\[0.3cm]
on the DEAP Dataset}\\[2cm]

{\large\bfseries Part I: Machine Learning Pipeline}\\[2cm]

\vfill

{\large
\begin{tabular}{ll}
\textbf{Submitted By:}  & [Student Names] \\
\textbf{Supervisor(s):} & [Supervisor Names] \\
\textbf{Submission Date:} & [Date] \\
\end{tabular}
}

\end{titlepage}

% ---- DECLARATION ----
\newpage
\thispagestyle{empty}
\chapter*{Declaration}
\addcontentsline{toc}{chapter}{Declaration}

We hereby certify that the work presented in this report titled
\textit{``EEG-Based Emotion Recognition Using Classical Machine Learning on the DEAP Dataset''}
is our own original work. All sources of information have been duly cited and
acknowledged. This report has not been previously submitted for any other degree
or qualification.

\vspace{3cm}
\noindent
\begin{tabular}{p{6cm}p{6cm}}
Signature: \hrulefill & Date: \hrulefill \\
\end{tabular}

% ---- ABSTRACT ----
\newpage
\chapter*{Abstract}
\addcontentsline{toc}{chapter}{Abstract}

This report presents the machine learning component of a two-part graduation
project aimed at building a real-time EEG-based emotion recognition system.
We design, implement, and systematically evaluate a classical machine learning
pipeline for arousal and valence classification using the DEAP
(Database for Emotion Analysis using Physiological signals) dataset.

The pipeline follows methodologies from high-accuracy published research and
encompasses signal preprocessing (bandpass filtering at 4--45~Hz, 50~Hz notch
filtering, and baseline correction), multi-domain feature extraction
(Differential Entropy, Higuchi Fractal Dimension, statistical moments, Hjorth
parameters, band power, and connectivity features), and classification using
gradient boosting and support vector machines with nested cross-validation.

Through a series of over 20 systematic experiments spanning window size
optimization, feature ablation, band selection, feature selection methods, and
model comparison, we achieved a best subject-dependent accuracy of
\textbf{66.72\% for valence} and \textbf{66.25\% for arousal} using XGBoost
with Differential Entropy and Higuchi Fractal Dimension features under rigorous
grouped cross-validation that prevents data leakage. We further demonstrate
that naive evaluation protocols that ignore trial boundaries can inflate
accuracy to over 80\%, highlighting the critical importance of proper
cross-validation in EEG research.

The trained models and feature extraction pipeline are designed to integrate
with the real-time mobile application presented in Part~II of this project.

\vspace{0.5cm}
\noindent\textbf{Keywords:} EEG, emotion recognition, DEAP dataset, machine
learning, differential entropy, Higuchi fractal dimension, XGBoost, SVM,
affective computing, brain-computer interface.

% ---- ACKNOWLEDGMENT ----
\newpage
\chapter*{Acknowledgment}
\addcontentsline{toc}{chapter}{Acknowledgment}

[We would like to express our gratitude to our supervisor(s) for their guidance
and support throughout this project. We also thank the creators of the DEAP
dataset for making their data publicly available for research purposes.]

% ---- TABLE OF CONTENTS ----
\newpage
\tableofcontents

% ---- LIST OF FIGURES ----
\newpage
\listoffigures
\addcontentsline{toc}{chapter}{List of Figures}

% ---- LIST OF TABLES ----
\newpage
\listoftables
\addcontentsline{toc}{chapter}{List of Tables}

% ---- LIST OF ABBREVIATIONS ----
\newpage
\chapter*{List of Abbreviations}
\addcontentsline{toc}{chapter}{List of Abbreviations}

\begin{longtable}{@{}ll@{}}
\toprule
\textbf{Abbreviation} & \textbf{Description} \\
\midrule
\endhead
BCI   & Brain-Computer Interface \\
CV    & Cross-Validation \\
DE    & Differential Entropy \\
DEAP  & Database for Emotion Analysis using Physiological signals \\
EEG   & Electroencephalography \\
FFT   & Fast Fourier Transform \\
HFD   & Higuchi Fractal Dimension \\
HGB   & Histogram-based Gradient Boosting \\
IIR   & Infinite Impulse Response \\
LGBM  & Light Gradient Boosting Machine \\
MI    & Mutual Information \\
ML    & Machine Learning \\
MLP   & Multi-Layer Perceptron \\
PCA   & Principal Component Analysis \\
PCC   & Pearson Correlation Coefficient \\
PLV   & Phase-Locking Value \\
PSD   & Power Spectral Density \\
RF    & Random Forest \\
ROC-AUC & Receiver Operating Characteristic -- Area Under Curve \\
SVM   & Support Vector Machine \\
XGB   & Extreme Gradient Boosting (XGBoost) \\
\bottomrule
\end{longtable}


% #############################################################################
% CHAPTER 1: INTRODUCTION
% #############################################################################
\chapter{Introduction}
\label{ch:introduction}

\section{Motivation}
\label{sec:motivation}

Emotions play a fundamental role in human cognition, decision-making, and
social interaction. The ability to automatically recognize emotional states
from physiological signals opens significant opportunities in healthcare,
human-computer interaction, adaptive learning systems, and mental health
monitoring. Among the various modalities available for emotion sensing,
electroencephalography (EEG) offers a direct window into brain activity,
capturing neural correlates of affective states with high temporal resolution.

Unlike facial expression or speech-based approaches, EEG signals are difficult
to consciously suppress or fake, making them a more reliable indicator of
genuine emotional experience. However, EEG-based emotion recognition presents
substantial challenges: signals are noisy, non-stationary, and exhibit high
inter-subject variability. Building robust classifiers that generalize across
individuals remains an open research problem.

The DEAP dataset~\cite{koelstra2012deap}, one of the most widely used
benchmarks in affective computing, provides multi-channel EEG recordings from
32 participants watching 40 music video stimuli, with self-reported arousal
and valence ratings. Despite being extensively studied, reported accuracies
in the literature vary widely---from below 60\% to above 90\%---with
discrepancies often attributable to differences in evaluation protocols rather
than genuine algorithmic improvements.

This project is motivated by two goals: (1)~building a rigorous, reproducible
ML pipeline for EEG emotion recognition that achieves competitive accuracy
under proper evaluation, and (2)~developing that pipeline into a deployable
system integrated with a real-time mobile application for practical use.

\section{Problem Definition}
\label{sec:problem}

Given multi-channel EEG recordings from the DEAP dataset, the task is to
classify emotional states along two affective dimensions:

\begin{itemize}[noitemsep]
    \item \textbf{Valence}: the pleasantness of an emotion (positive vs.\ negative).
    \item \textbf{Arousal}: the intensity or activation level of an emotion (calm vs.\ excited).
\end{itemize}

Each dimension is treated as a binary classification problem using a threshold
on the self-reported ratings (scale 1--9). The pipeline must address several
sub-problems:

\begin{enumerate}[noitemsep]
    \item \textbf{Signal preprocessing}: Removing noise and artifacts from raw EEG.
    \item \textbf{Feature extraction}: Computing discriminative representations
          from time-domain, frequency-domain, and spatial perspectives.
    \item \textbf{Model selection and evaluation}: Identifying the best classifier
          and hyperparameters using cross-validation that avoids data leakage.
    \item \textbf{Deployability}: Producing a model suitable for integration
          with a real-time mobile application.
\end{enumerate}

\section{Objectives}
\label{sec:objectives}

The objectives of this project are:

\begin{enumerate}[noitemsep]
    \item Implement a modular, end-to-end ML pipeline for EEG-based emotion
          recognition following published methodologies from high-accuracy papers.
    \item Systematically evaluate the impact of preprocessing parameters,
          feature types, window sizes, feature selection methods, and classifiers
          on recognition accuracy.
    \item Achieve competitive subject-dependent classification accuracy on the
          DEAP dataset under rigorous, leakage-free evaluation.
    \item Demonstrate the effect of data leakage on reported results to
          contextualize discrepancies in published literature.
    \item Provide a trained model and feature extraction pipeline suitable for
          real-time deployment in the mobile application (Part~II).
\end{enumerate}

\section{Contributions}
\label{sec:contributions}

The main contributions of this work are:

\begin{enumerate}[noitemsep]
    \item A \textbf{modular, open-source pipeline} with configurable preprocessing,
          17 feature extraction methods, 10 classifiers, and multiple evaluation
          protocols, with built-in caching and checkpoint/resume capabilities.
    \item A \textbf{systematic experimental study} comprising over 20 experiments
          with ablation across window sizes, feature sets, frequency bands,
          feature selection methods, and models.
    \item A \textbf{rigorous evaluation methodology} using grouped
          cross-validation (GroupKFold by trial) that prevents data leakage
          between overlapping windows from the same trial.
    \item A \textbf{quantitative demonstration of data leakage effects},
          showing that improper evaluation inflates accuracy by approximately
          15--20 percentage points.
    \item \textbf{Reproducible results} with checkpointing, caching, and
          configuration serialization for full experimental traceability.
\end{enumerate}

\section{Report Organization}
\label{sec:organization}

This report is organized as follows:

\begin{itemize}[noitemsep]
    \item \textbf{Chapter~\ref{ch:background}}: Reviews theoretical foundations
          of EEG signal processing, feature extraction techniques, and machine
          learning classifiers. Presents a literature review of DEAP-based
          emotion recognition work.
    \item \textbf{Chapter~\ref{ch:architecture}}: Describes the system
          architecture, data flow, and design decisions of the pipeline.
    \item \textbf{Chapter~\ref{ch:experiments}}: Details all experiments
          conducted, including configurations, results, and analysis.
    \item \textbf{Chapter~\ref{ch:conclusion}}: Summarizes findings and
          discusses integration with the mobile application (Part~II).
\end{itemize}

\noindent\textit{Note: This report covers Part~I (Machine Learning Pipeline)
of the graduation project. Part~II (Real-Time Mobile Application) will be
appended as additional chapters.}


% #############################################################################
% CHAPTER 2: THEORETICAL BACKGROUND
% #############################################################################
\chapter{Theoretical Background}
\label{ch:background}

\section{Background}
\label{sec:bg}

\subsection{Emotion Models}

Emotions can be represented using discrete categories (e.g., happiness, sadness,
anger) or continuous dimensional models. The most widely adopted dimensional
model is Russell's \textit{circumplex model of affect}~\cite{russell1980}, which
describes emotions along two orthogonal axes:

\begin{itemize}[noitemsep]
    \item \textbf{Valence}: Ranges from unpleasant (negative) to pleasant
          (positive).
    \item \textbf{Arousal}: Ranges from calm (low activation) to excited
          (high activation).
\end{itemize}

This two-dimensional representation allows any emotion to be mapped as a point
in the valence-arousal space. For example, ``excitement'' corresponds to high
valence and high arousal, while ``sadness'' corresponds to low valence and low
arousal. The DEAP dataset adopts this model, providing self-reported valence
and arousal ratings on a continuous 1--9 scale.

\subsection{Electroencephalography (EEG)}

Electroencephalography records electrical activity of the brain through
electrodes placed on the scalp. The DEAP dataset uses 32 EEG channels arranged
according to the international 10--20 system (channels: Fp1, AF3, F7, F3, FC1,
FC5, T7, C3, CP1, CP5, P7, P3, Pz, PO3, O1, Oz, O2, PO4, P4, P8, CP6, CP2,
C4, T8, FC6, FC2, F4, F8, AF4, Fp2, Fz, Cz), sampled at 128~Hz and
downsampled from the original 512~Hz recordings.

EEG signals are commonly analyzed in terms of frequency bands, each associated
with different cognitive and affective states:

\begin{table}[H]
\centering
\caption{EEG frequency bands and their functional associations.}
\label{tab:eeg_bands}
\begin{tabular}{llp{7cm}}
\toprule
\textbf{Band} & \textbf{Range (Hz)} & \textbf{Functional Association} \\
\midrule
Theta  & 4--8   & Memory processing, emotional regulation, drowsiness \\
Alpha  & 8--13  & Relaxation, idling, inhibition of task-irrelevant regions \\
Beta   & 13--30 & Active thinking, focus, anxiety, motor planning \\
Gamma  & 30--45 & Higher cognitive functions, cross-modal sensory processing \\
\bottomrule
\end{tabular}
\end{table}

\subsection{The DEAP Dataset}

The DEAP dataset~\cite{koelstra2012deap} is a benchmark for emotion
recognition from physiological signals. Its key characteristics are:

\begin{itemize}[noitemsep]
    \item \textbf{Participants}: 32 subjects.
    \item \textbf{Stimuli}: 40 one-minute music video excerpts selected to
          elicit a range of emotional responses.
    \item \textbf{EEG channels}: 32 channels at 128~Hz sampling rate.
    \item \textbf{Trial structure}: Each trial consists of a 3-second
          pre-stimulus baseline followed by 60 seconds of stimulus viewing.
    \item \textbf{Self-assessment}: After each trial, participants rated their
          experience on 1--9 scales for valence, arousal, dominance, and liking.
    \item \textbf{Data format}: Each subject's data is stored as a
          \texttt{.mat} file containing EEG data with shape
          $[40 \times 40 \times 8064]$ (trials $\times$ channels $\times$
          samples) and labels with shape $[40 \times 4]$.
\end{itemize}

The 8064 samples per trial correspond to 63~seconds at 128~Hz (3~s baseline +
60~s stimulus). After baseline removal, 60~seconds (7680~samples) of stimulus
data remain per trial.

\subsection{Classical Machine Learning for EEG}

Classical ML approaches for EEG classification typically follow a pipeline of:
(1)~preprocessing, (2)~feature extraction, (3)~optional feature
selection/reduction, and (4)~classification. Unlike deep learning approaches,
classical methods rely on hand-crafted features derived from domain knowledge
of EEG signal characteristics. This approach has several advantages for
EEG-based emotion recognition:

\begin{itemize}[noitemsep]
    \item \textbf{Interpretability}: Features have clear neurophysiological
          meaning (e.g., alpha power relates to relaxation).
    \item \textbf{Sample efficiency}: Works with the relatively small sample
          sizes typical of EEG datasets (32 subjects $\times$ 40 trials).
    \item \textbf{Computational efficiency}: Suitable for real-time deployment
          on mobile devices.
    \item \textbf{Competitive accuracy}: State-of-the-art classical methods
          achieve comparable or superior results to deep learning on DEAP.
\end{itemize}


\section{Literature Review}
\label{sec:literature}

EEG-based emotion recognition on the DEAP dataset has been extensively studied.
We review the key papers that informed our pipeline design, focusing on those
reporting the highest accuracies with classical ML methods.

\subsection{High-Accuracy Approaches}

\textbf{Paper~13} (Differential Entropy and Fractal Dimension features)
proposed using histogram-based Differential Entropy (DE) and Higuchi Fractal
Dimension (HFD) as features, combined with XGBoost classification. Their
preprocessing included a 4th-order Butterworth bandpass filter (4--45~Hz) and
a 50~Hz notch filter for power line interference removal. Using 3-second
non-overlapping windows, they reported \textbf{89\% valence and 88\% arousal}
accuracy. This paper formed the primary basis for our pipeline design.

\textbf{Paper~14} (Statistical and Hjorth features) achieved \textbf{91\%
valence and 89\% arousal} using statistical features (mean, standard deviation,
skewness, kurtosis) and Hjorth parameters (activity, mobility, complexity)
combined with XGBoost. Their high accuracies motivated the inclusion of these
feature types in our pipeline.

\textbf{Paper~11} (Band power with feature selection) used band power features
across the standard EEG frequency bands with Mutual Information-based feature
selection, reporting \textbf{64--89\%} accuracy depending on the evaluation
protocol and classifier.

\subsection{Evaluation Protocol Concerns}

A critical observation from the literature is the wide variation in reported
accuracies (55\%--95\%) for the same dataset, which is largely attributable to
differences in evaluation protocols:

\begin{itemize}[noitemsep]
    \item \textbf{Window-level vs.\ trial-level evaluation}: When 60-second
          trials are segmented into 3-second windows, adjacent windows from the
          same trial are highly correlated. Random train/test splits that mix
          windows from the same trial into both sets create data leakage.
    \item \textbf{Subject-dependent vs.\ cross-subject}: Subject-dependent
          evaluation (training and testing on the same subject) yields higher
          accuracy than cross-subject generalization.
    \item \textbf{Pooled vs.\ per-subject evaluation}: Pooling all subjects'
          data without accounting for subject identity inflates accuracy.
\end{itemize}

Rigorous evaluation requires \textbf{grouped cross-validation} (GroupKFold)
where all windows from the same trial appear exclusively in either the training
or test set, never both.


\section{State of the Art}
\label{sec:sota}

Table~\ref{tab:sota} summarizes representative results from the literature on
the DEAP dataset using classical ML methods.

\begin{table}[H]
\centering
\caption{Summary of representative DEAP emotion recognition results.}
\label{tab:sota}
\begin{tabular}{p{3.5cm}lllc}
\toprule
\textbf{Method} & \textbf{Features} & \textbf{Classifier} & \textbf{Val / Aro} & \textbf{CV Type} \\
\midrule
Paper 13 & DE, HFD & XGBoost & 89\% / 88\% & Not grouped \\
Paper 14 & Statistical, Hjorth & XGBoost & 91\% / 89\% & Not grouped \\
Paper 11 & Band power & SVM & 64--89\% & Varies \\
\midrule
\textbf{Ours (rigorous)} & \textbf{DE, HFD} & \textbf{XGBoost} & \textbf{66.7\% / 66.3\%} & \textbf{Grouped} \\
Ours (leaky, ref.) & DE, BP, Asym & XGBoost & 80.9\% / 81.9\% & Not grouped \\
\bottomrule
\end{tabular}
\end{table}

The gap between our rigorous results and the literature values is consistent
with the data leakage hypothesis: when we deliberately replicate the leaky
evaluation protocol, our accuracy rises to a comparable $\sim$81\%.


% #############################################################################
% CHAPTER 3: SYSTEM ARCHITECTURE AND DESIGN
% #############################################################################
\chapter{System Architecture and Design}
\label{ch:architecture}

\section{Introduction}

This chapter presents the architecture and design of the EEG emotion
recognition pipeline. The system is implemented in Python as a modular package
(\texttt{deap\_emotion}) with clear separation of concerns across preprocessing,
feature extraction, labeling, model building, and evaluation components.

\section{System Overview}

The pipeline follows the standard EEG processing workflow illustrated in
Figure~\ref{fig:pipeline}.

\begin{figure}[H]
\centering
\fbox{\parbox{0.9\textwidth}{\centering
\textbf{Pipeline Data Flow}\\[0.5cm]
DEAP Files (\texttt{archive/*.mat})\\
$\downarrow$\\
\texttt{data.py}: Load \& Parse $\rightarrow$ \texttt{SubjectData} (trials $\times$ channels $\times$ samples)\\
$\downarrow$\\
\texttt{preprocess.py}: Bandpass 4--45~Hz $\rightarrow$ Notch 50~Hz $\rightarrow$ Baseline Correction\\
$\downarrow$\\
\texttt{features.py}: Segmentation (3s windows) $\rightarrow$ Feature Extraction\\
$\downarrow$\\
\texttt{labels.py}: Build Binary Labels (threshold-based)\\
$\downarrow$\\
\texttt{model\_registry.py}: Build Classifier Pipeline (Scaler $\rightarrow$ Selection $\rightarrow$ Classifier)\\
$\downarrow$\\
\texttt{experiment.py}: Nested CV (Outer: Repeated Group Splits, Inner: Grid Search)\\
$\downarrow$\\
Reports: JSON + CSV + TXT summaries
}}
\caption{End-to-end data flow of the emotion recognition pipeline.}
\label{fig:pipeline}
\end{figure}

\section{Data Shapes and Conventions}

Throughout the pipeline, the following tensor conventions are maintained:

\begin{table}[H]
\centering
\caption{Data shape conventions used in the pipeline.}
\label{tab:shapes}
\begin{tabular}{lll}
\toprule
\textbf{Data} & \textbf{Shape} & \textbf{Typical Dimensions} \\
\midrule
Raw EEG         & [trials, channels, samples] & [40, 32, 8064] \\
Post-baseline   & [trials, channels, samples] & [40, 32, 7680] \\
Windowed EEG    & [segments, channels, window] & [800, 32, 384] \\
Feature matrix  & [samples, features]          & [800, 64] \\
Labels / Groups & [samples]                    & [800] \\
\bottomrule
\end{tabular}
\end{table}

\noindent where 800 segments arise from 40 trials $\times$ 20 windows/trial
(60s $\div$ 3s = 20) per subject.


\section{Preprocessing Module}
\label{sec:preprocess}

The preprocessing pipeline, implemented in \texttt{preprocess.py}, follows the
methodology of Paper~13 and consists of three stages:

\subsection{Bandpass Filtering}

A 4th-order Butterworth bandpass filter with cutoff frequencies at 4~Hz and
45~Hz is applied using zero-phase filtering (\texttt{scipy.signal.filtfilt}).
The filter design uses the bilinear transform:

\begin{equation}
    H(s) = \frac{1}{\prod_{k=1}^{N}(s - p_k)}
    \label{eq:butterworth}
\end{equation}

\noindent where $N=4$ is the filter order and $p_k$ are the Butterworth poles.
Zero-phase filtering applies the filter forward and backward, eliminating phase
distortion while doubling the effective order. The 4~Hz lower cutoff removes
low-frequency drift and the delta band (which is dominated by eye movement
artifacts), while the 45~Hz upper cutoff removes high-frequency muscle noise.

\subsection{Notch Filtering}

A 50~Hz IIR notch filter (quality factor $Q=30$) removes power line
interference using \texttt{scipy.signal.iirnotch}. The narrow bandwidth
($50 \pm 1.67$~Hz) ensures minimal distortion to the gamma band signal.

\subsection{Baseline Correction}

The first 3~seconds (384~samples at 128~Hz) of each trial serve as the
pre-stimulus baseline. The per-channel mean of this baseline period is
subtracted from the entire trial, after which the baseline samples are
discarded, leaving 60~seconds of stimulus data:

\begin{equation}
    x_{\text{corrected}}[c, t] = x[c, t] - \frac{1}{T_b}\sum_{t=1}^{T_b} x[c, t]
    \quad \text{for } t > T_b
\end{equation}

\noindent where $T_b = 384$ is the number of baseline samples and $c$ indexes
the channel.


\section{Feature Extraction Module}
\label{sec:features}

The feature extraction module (\texttt{features.py}) supports 12 feature types
organized into four categories. All features are computed per-channel and then
flattened into a single feature vector per window.

\subsection{Frequency-Domain Features}

\textbf{Band Power} (\texttt{bandpower}): The average power spectral density
within each frequency band, computed via FFT:

\begin{equation}
    P_b = \frac{1}{|F_b|}\sum_{f \in F_b} |X(f)|^2
\end{equation}

\noindent where $F_b$ is the set of frequency bins falling within band $b$.
This produces 4 features per channel (theta, alpha, beta, gamma) = 128
features for 32 channels.

\textbf{Differential Entropy} (\texttt{de\_histogram}): Computed using the
histogram-based method from Paper~13:

\begin{equation}
    h(X) = -\sum_{i=1}^{B} p(x_i) \cdot \log p(x_i) \cdot \Delta x
    \label{eq:de}
\end{equation}

\noindent where $B=100$ histogram bins, $p(x_i)$ is the probability density
estimate, and $\Delta x$ is the bin width. This produces 1 feature per channel
= 32 features.

\textbf{Relative Band Power} (\texttt{rel\_bandpower}): Band power normalized
by total power: $P_{b,\text{rel}} = P_b / \sum_b P_b$.

\textbf{PSD Bins} (\texttt{psd}): Fine-grained power spectral density in
2~Hz-wide bins from 1--45~Hz.

\subsection{Time-Domain Features}

\textbf{Statistical Features} (\texttt{statistical}, Paper~14): Per-channel
mean ($\mu$), standard deviation ($\sigma$), skewness ($\gamma_1$), and
kurtosis ($\gamma_2$). Produces 4 features per channel = 128 features.

\textbf{Hjorth Parameters} (\texttt{hjorth}, Paper~14): Three descriptors
per channel capturing signal shape:

\begin{align}
    \text{Activity} &= \text{Var}(x) \\
    \text{Mobility} &= \sqrt{\frac{\text{Var}(x')}{\text{Var}(x)}} \\
    \text{Complexity} &= \frac{\text{Mobility}(x')}{\text{Mobility}(x)}
\end{align}

\noindent Produces 3 features per channel = 96 features.

\textbf{Higuchi Fractal Dimension} (\texttt{hfd}, Paper~13): Quantifies signal
complexity by measuring the fractal dimension through curve length analysis at
multiple scales ($k = 1, \ldots, k_{\max}=10$):

\begin{equation}
    L(k) = \frac{1}{k} \sum_{m=1}^{k} \frac{(n-1) \sum_{i=1}^{\lfloor(n-m)/k\rfloor}
    |x(m+ik) - x(m+(i-1)k)|}{k \cdot \lfloor(n-m)/k\rfloor \cdot k}
\end{equation}

The fractal dimension $D$ is the slope of $\log L(k)$ vs.\ $\log(1/k)$.
Produces 1 feature per channel = 32 features.

\subsection{Spatial/Connectivity Features}

\textbf{Asymmetry} (\texttt{asymmetry}): Differential and ratio features
between 7 symmetric electrode pairs (Fp1--Fp2, F3--F4, F7--F8, C3--C4,
T7--T8, P3--P4, O1--O2), capturing hemispheric lateralization of emotional
processing.

\textbf{Pearson Correlation Coefficient} (\texttt{connectivity\_pcc}):
Summary statistics (mean, std, mean absolute) of pairwise channel correlations
within each frequency band, capturing functional connectivity.

\textbf{Phase-Locking Value} (\texttt{connectivity\_plv}): Summary statistics
of pairwise phase synchronization computed via the Hilbert transform, measuring
phase coherence between channels.

\subsection{Feature Dimensionality Summary}

\begin{table}[H]
\centering
\caption{Feature dimensionality for 32 EEG channels and 4 frequency bands.}
\label{tab:feature_dims}
\begin{tabular}{llr}
\toprule
\textbf{Feature Type} & \textbf{Per Channel} & \textbf{Total Features} \\
\midrule
Band Power         & 4 (per band)   & 128 \\
Differential Entropy (histogram) & 1 & 32 \\
Statistical        & 4              & 128 \\
Hjorth Parameters  & 3              & 96 \\
Higuchi FD         & 1              & 32 \\
Asymmetry          & 4 (per pair)   & 56 \\
PCC Connectivity   & 3 (per band)   & 12 \\
PLV Connectivity   & 3 (per band)   & 12 \\
\bottomrule
\end{tabular}
\end{table}


\section{Labeling Module}
\label{sec:labels}

The labeling module (\texttt{labels.py}) converts continuous self-report ratings
into classification targets. Seven modes are supported:

\begin{itemize}[noitemsep]
    \item \textbf{Binary}: $y = 1$ if rating $>$ threshold, else $y = 0$.
          Default threshold = 4.5 (Paper~13).
    \item \textbf{Three-class}: Low / neutral / high based on configurable bins
          (default: $[3.0, 6.0]$).
    \item \textbf{Quartile}: Four classes based on rating quartiles.
    \item \textbf{Subject mean/median}: Per-subject adaptive threshold using
          the individual's mean or median rating.
    \item \textbf{Subject tertile}: Three classes based on per-subject tertiles.
    \item \textbf{Regression}: Raw continuous ratings (for regression models).
\end{itemize}

The primary experiments use \textbf{binary} labeling with a threshold of 4.5.


\section{Model Registry}
\label{sec:models}

The model registry (\texttt{model\_registry.py}) defines 10 classifiers and 4
regression models as \texttt{ModelSpec} dataclasses, each specifying default
parameters, hyperparameter search grids, and pipeline requirements.

\begin{table}[H]
\centering
\caption{Available classifiers and their configurations.}
\label{tab:models}
\begin{tabular}{llcp{4.5cm}}
\toprule
\textbf{Model} & \textbf{Algorithm} & \textbf{Scaler} & \textbf{Key Hyperparameters} \\
\midrule
\texttt{xgb} & XGBoost & No & n\_estimators, max\_depth, learning\_rate, subsample \\
\texttt{rbf\_svm} & RBF-SVM & Yes & C, gamma, class\_weight \\
\texttt{rf} & Random Forest & No & n\_estimators, max\_depth \\
\texttt{hgb} & Hist. Gradient Boost & No & max\_depth, learning\_rate, max\_iter \\
\texttt{logreg} & Logistic Regression & Yes & C \\
\texttt{linear\_svm} & Linear SVM & Yes & C \\
\texttt{mlp} & Multi-Layer Perceptron & Yes & hidden\_layers, alpha \\
\texttt{lgbm} & LightGBM & No & n\_estimators, learning\_rate, num\_leaves \\
\texttt{catboost} & CatBoost & No & iterations, learning\_rate, depth \\
\texttt{soft\_vote} & Soft Voting Ensemble & --- & Combines: RBF-SVM, LogReg, HGB \\
\bottomrule
\end{tabular}
\end{table}

Each model is wrapped in a \texttt{sklearn.pipeline.Pipeline} that optionally
includes: (1)~StandardScaler or RobustScaler, (2)~feature selection via Mutual
Information (\texttt{SelectKBest}) or L1 regularization (\texttt{SelectFromModel}),
(3)~PCA dimensionality reduction, and (4)~the classifier.


\section{Evaluation Module}
\label{sec:eval}

\subsection{Subject-Dependent Evaluation}

The primary evaluation protocol is \textbf{subject-dependent} with repeated
grouped random splits:

\begin{enumerate}[noitemsep]
    \item For each subject, split the 40 trials into 80\% training (32 trials)
          and 20\% testing (8 trials) using random permutation.
    \item All windows from a given trial belong exclusively to either the
          training or test set (GroupKFold by trial ID), preventing leakage.
    \item Perform inner cross-validation (GroupKFold) on the training set for
          hyperparameter tuning via grid search.
    \item Evaluate on the held-out test trials.
    \item Repeat for 5 outer splits with different random seeds.
    \item Report trial-level accuracy using majority voting across windows.
\end{enumerate}

\subsection{Cross-Subject Evaluation}

Cross-subject evaluation uses Leave-One-Group-Out, where each subject serves
as one fold:

\begin{enumerate}[noitemsep]
    \item Leave one subject out for testing.
    \item Train on all remaining subjects.
    \item Repeat for all 32 subjects.
\end{enumerate}

\subsection{Trial-Level Aggregation}

Since each 60-second trial is segmented into multiple 3-second windows, the
final prediction for a trial is obtained by \textbf{majority voting}: the class
predicted most frequently across all windows of that trial becomes the trial's
prediction. An alternative is \textbf{mean probability} aggregation, where
predicted probabilities are averaged before thresholding.


\section{Caching and Checkpointing}
\label{sec:caching}

\textbf{Feature caching}: Extracted features are cached to disk using a SHA-256
hash of all feature-relevant configuration parameters (channels, bands, window
size, feature set, preprocessing flags). This ensures cache invalidation when
any parameter changes while avoiding redundant computation.

\textbf{Experiment checkpointing}: After each fold evaluation, results are
appended to a checkpoint CSV file. On restart, completed folds are loaded from
the checkpoint and skipped, enabling experiments to be paused and resumed
without loss of progress. A live summary file is updated after each fold,
providing real-time monitoring of running experiments.


\section{Software Architecture}

The codebase is organized as a Python package with the following structure:

\begin{table}[H]
\centering
\caption{Source code modules and their responsibilities.}
\label{tab:modules}
\begin{tabular}{lp{9cm}}
\toprule
\textbf{Module} & \textbf{Responsibility} \\
\midrule
\texttt{config.py}          & Frozen dataclass with all pipeline parameters \\
\texttt{data.py}            & Dataset loading (MAT/NPZ/DAT formats) \\
\texttt{preprocess.py}      & Bandpass, notch, baseline, segmentation \\
\texttt{features.py}        & 12 feature extraction methods \\
\texttt{labels.py}          & 7 labeling modes \\
\texttt{model\_registry.py} & 14 model definitions with hyperparameter grids \\
\texttt{eval.py}            & Subject-dependent \& cross-subject evaluation \\
\texttt{experiment.py}      & Grid search, nested CV, checkpointing \\
\texttt{metrics.py}         & 20+ classification and regression metrics \\
\texttt{splits.py}          & GroupKFold and LeaveOneGroupOut splitters \\
\texttt{riemann.py}         & Riemannian tangent space transform \\
\texttt{cli\_experiment.py} & Command-line interface \\
\bottomrule
\end{tabular}
\end{table}

Key dependencies include NumPy, SciPy (signal processing), scikit-learn
(ML framework), XGBoost, and h5py (DEAP file loading).


% #############################################################################
% CHAPTER 4: EXPERIMENTS AND RESULTS
% #############################################################################
\chapter{Experiments and Results}
\label{ch:experiments}

This chapter presents the systematic experiments conducted to optimize the
emotion recognition pipeline. Over 20 experiment configurations were evaluated,
organized into four experimental series: (1)~Paper~13 replication,
(2)~window and feature optimization (roadmap), (3)~representation learning,
and (4)~data leakage analysis.

All experiments use the full DEAP dataset (32 subjects, 40 trials each) unless
otherwise noted. The primary metric is \textbf{trial-level accuracy} obtained
by majority voting across windows within each trial.


\section{Experiment 1: Paper 13 Replication}
\label{sec:exp_paper13}

\subsection{Configuration}

This experiment replicates the methodology of Paper~13 as closely as possible:

\begin{table}[H]
\centering
\caption{Paper 13 replication configuration.}
\label{tab:paper13_config}
\begin{tabular}{ll}
\toprule
\textbf{Parameter} & \textbf{Value} \\
\midrule
Preprocessing & Bandpass 4--45~Hz (Butterworth 4th order), Notch 50~Hz \\
Baseline      & 3~s subtraction \\
Window / Step & 3.0~s / 3.0~s (no overlap) \\
Features      & DE (histogram) + HFD \\
Feature dim.  & 64 (32 DE + 32 HFD) \\
Classifier    & XGBoost (default params) \\
Label         & Binary, threshold = 4.5 \\
CV            & 5-fold GroupKFold by trial (within-subject) \\
Subjects      & All 32 \\
\bottomrule
\end{tabular}
\end{table}

\subsection{Results --- Default Parameters (Fast)}

With default XGBoost parameters (no hyperparameter tuning):

\begin{table}[H]
\centering
\caption{Paper 13 fast replication results (default XGBoost, 32 subjects).}
\label{tab:paper13_fast}
\begin{tabular}{lcc}
\toprule
\textbf{Dimension} & \textbf{Accuracy} & \textbf{Std Dev} \\
\midrule
Valence & 61.4\% & $\pm$ 7.0\% \\
Arousal & 62.9\% & $\pm$ 11.5\% \\
\bottomrule
\end{tabular}
\end{table}

\subsection{Results --- With Grid Search}

With nested hyperparameter tuning (5 outer repeats, 4 inner folds):

\begin{table}[H]
\centering
\caption{Paper 13 replication with grid search (160 folds, 32 subjects).}
\label{tab:paper13_grid}
\begin{tabular}{lcccc}
\toprule
\textbf{Dimension} & \textbf{Accuracy} & \textbf{Balanced Acc.} & \textbf{F1 (Macro)} & \textbf{Std Dev} \\
\midrule
Valence & 66.72\% & 58.43\% & 53.25\% & $\pm$ 8.94\% \\
Arousal & 66.25\% & 56.81\% & 51.31\% & $\pm$ 13.30\% \\
\bottomrule
\end{tabular}
\end{table}

Grid search improved accuracy by approximately 5 percentage points over
default parameters, demonstrating the importance of hyperparameter tuning.
This is the \textbf{best result} achieved under rigorous evaluation.

\subsection{Analysis}

The gap between our result (66.7\%) and Paper~13's reported result (89\%) is
substantial. We attribute this primarily to differences in evaluation
protocol: Paper~13 does not clearly specify grouped cross-validation, and
our leakage analysis (Section~\ref{sec:leakage}) shows that non-grouped
evaluation can inflate accuracy by $\sim$15--20 percentage points.

The balanced accuracy (58.4\%) is notably lower than raw accuracy (66.7\%),
indicating class imbalance in the binary split at threshold 4.5. The F1
macro score (53.3\%) reflects difficulty in the minority class.


\section{Experiment 2: Window Size Optimization}
\label{sec:exp_window}

We systematically varied the window size to find the optimal temporal
resolution for feature extraction.

\subsection{Roadmap Phased --- Phase 1 Window Sweep}

Using RBF-SVM with DE features (2 outer repeats, 3 inner folds, subset of
subjects for quick screening):

\begin{table}[H]
\centering
\caption{Window size sweep results (RBF-SVM, DE features, Roadmap Phased).}
\label{tab:window_sweep_phased}
\begin{tabular}{llcc}
\toprule
\textbf{Config} & \textbf{Window / Step} & \textbf{Valence} & \textbf{Arousal} \\
\midrule
w2\_s2 & 2.0s / 2.0s & 60.00\% & 53.12\% \\
w2\_s1 & 2.0s / 1.0s & 60.00\% & 55.63\% \\
w4\_s4 & 4.0s / 4.0s & 57.50\% & 54.37\% \\
w4\_s2 & 4.0s / 2.0s & 53.75\% & 55.00\% \\
w8\_s8 & 8.0s / 8.0s & 60.62\% & 50.62\% \\
w8\_s4 & 8.0s / 4.0s & 58.75\% & 53.12\% \\
\bottomrule
\end{tabular}
\end{table}

\subsection{Roadmap Ceiling --- Full Window Sweep}

Extended sweep on all 32 subjects with RBF-SVM (3 outer repeats):

\begin{table}[H]
\centering
\caption{Window size sweep results (RBF-SVM, DE features, all 32 subjects).}
\label{tab:window_sweep_full}
\begin{tabular}{llcc}
\toprule
\textbf{Config} & \textbf{Window / Step} & \textbf{Valence} & \textbf{Arousal} \\
\midrule
full\_w2\_s1  & 2.0s / 1.0s  & 57.03\% & 59.51\% \\
full\_w4\_s2  & 4.0s / 2.0s  & 57.03\% & 59.24\% \\
full\_w10\_s5 & 10.0s / 5.0s & 58.59\% & 58.72\% \\
full\_w15\_s7.5 & 15.0s / 7.5s & 56.38\% & 59.64\% \\
\bottomrule
\end{tabular}
\end{table}

\subsection{Analysis}

Window size had a moderate effect on accuracy. The 3-second window (Paper~13's
choice) performed well for XGBoost, while the 6-second window showed marginal
improvement for RBF-SVM. Very large windows (15~s) reduced performance, likely
due to fewer training samples and temporal blurring of emotional responses.
Overlapping windows (step $<$ window) did not yield significant improvement
under grouped CV, since windows from the same trial cannot appear in both
train and test sets.


\section{Experiment 3: Feature Ablation}
\label{sec:exp_features}

We evaluated different feature combinations to identify the most discriminative
representations.

\subsection{Roadmap Focus --- Feature Ablation (6s window, RBF-SVM)}

\begin{table}[H]
\centering
\caption{Feature ablation results (6s window, 3s step, RBF-SVM).}
\label{tab:feature_ablation}
\begin{tabular}{lcc}
\toprule
\textbf{Features} & \textbf{Valence} & \textbf{Arousal} \\
\midrule
DE only              & \textbf{61.88\%} & \textbf{57.50\%} \\
DE + Asymmetry       & 60.00\% & 58.75\% \\
DE + Band Power      & 56.25\% & 54.37\% \\
\bottomrule
\end{tabular}
\end{table}

\subsection{Roadmap Ceiling --- Representation Learning (10s window, RBF-SVM)}

\begin{table}[H]
\centering
\caption{Representation learning results (10s/5s window, RBF-SVM, 160 folds).}
\label{tab:repr_learning}
\begin{tabular}{lcc}
\toprule
\textbf{Features} & \textbf{Valence} & \textbf{Arousal} \\
\midrule
DE only              & 57.42\% & 57.42\% \\
DE + Asymmetry       & 57.66\% & 57.19\% \\
DE + PCC             & 57.81\% & 56.72\% \\
DE + PLV             & 58.36\% & 57.42\% \\
DE + PCC + PLV       & 58.13\% & \textbf{58.44\%} \\
DE + PLV + MI (k=100) & \textbf{59.06\%} & 57.73\% \\
DE + PLV + MI (k=200) & 58.20\% & 57.34\% \\
DE + PLV + PCA (0.95) & 58.05\% & 57.58\% \\
\bottomrule
\end{tabular}
\end{table}

\subsection{Analysis}

A consistent finding across experiments: \textbf{Differential Entropy alone
outperforms feature combinations}. Adding band power, asymmetry, or
connectivity features either maintains or reduces accuracy. This suggests
that DE captures the most emotionally relevant information, and additional
features introduce noise that degrades generalization. Connectivity features
(PLV, PCC) provide marginal improvement ($\sim$0.5--1\%) in some
configurations but increase dimensionality substantially.


\section{Experiment 4: Band Selection}
\label{sec:exp_bands}

We tested whether excluding the delta band (which is largely removed by the
4~Hz highpass filter) would affect results.

\begin{table}[H]
\centering
\caption{Band selection results (DE only, 6s/3s window, RBF-SVM).}
\label{tab:band_selection}
\begin{tabular}{lcc}
\toprule
\textbf{Band Set} & \textbf{Valence} & \textbf{Arousal} \\
\midrule
Default (theta, alpha, beta, gamma) & 61.88\% & 57.50\% \\
No delta (identical, since delta excluded by bandpass) & 61.88\% & 57.50\% \\
\bottomrule
\end{tabular}
\end{table}

The results confirm that the delta band contributes no information after
4~Hz highpass filtering, validating the four-band (theta, alpha, beta, gamma)
configuration.


\section{Experiment 5: Feature Selection}
\label{sec:exp_selection}

We compared feature selection methods to determine whether dimensionality
reduction improves generalization.

\begin{table}[H]
\centering
\caption{Feature selection comparison (DE only, 6s/3s, RBF-SVM).}
\label{tab:feature_selection}
\begin{tabular}{lcc}
\toprule
\textbf{Method} & \textbf{Valence} & \textbf{Arousal} \\
\midrule
None (all features) & \textbf{61.88\%} & \textbf{57.50\%} \\
MI (k=200)          & 61.88\% & 57.50\% \\
L1 (C=0.1)          & 60.62\% & 52.50\% \\
L1 (C=0.5)          & 61.88\% & 51.25\% \\
L1 (C=1.0)          & 61.88\% & 55.00\% \\
\bottomrule
\end{tabular}
\end{table}

Feature selection provided no benefit for DE-only features (32 dimensions),
as the feature space is already low-dimensional. MI selection with $k=200$
retained all features since $k > 32$. L1-based selection reduced arousal
accuracy, possibly by discarding informative channels.


\section{Experiment 6: Model Comparison}
\label{sec:exp_models}

The final roadmap experiment compared models using the best configuration
(DE only, 6s/3s window, 5 outer repeats, 4 inner folds):

\begin{table}[H]
\centering
\caption{Model comparison (DE only, 6s/3s window, all 32 subjects, 160 folds).}
\label{tab:model_comparison}
\begin{tabular}{lcccc}
\toprule
\textbf{Model} & \textbf{Valence} & \textbf{Arousal} & \textbf{Val Std} & \textbf{Aro Std} \\
\midrule
RBF-SVM     & 58.98\% & \textbf{58.67\%} & $\pm$ 13.1\% & $\pm$ 10.8\% \\
HGB         & \textbf{59.84\%} & 56.48\% & $\pm$ 11.6\% & $\pm$ 9.9\% \\
Soft Voting & 59.38\% & 57.42\% & $\pm$ 12.2\% & $\pm$ 10.9\% \\
\bottomrule
\end{tabular}
\end{table}

\noindent
Comparing with the Paper~13 configuration (XGBoost, DE+HFD, 3s window, grid
search): \textbf{66.72\% / 66.25\%} remains the best overall result, with
XGBoost outperforming both RBF-SVM and HGB when paired with the HFD feature.


\section{Experiment 7: Data Leakage Analysis}
\label{sec:leakage}

To explain the gap between our results and published accuracies, we
deliberately introduced data leakage by pooling all windows across subjects
and trials and using stratified random splits (ignoring trial boundaries).

\begin{table}[H]
\centering
\caption{Data leakage analysis: effect of evaluation protocol on accuracy.}
\label{tab:leakage}
\begin{tabular}{lcccc}
\toprule
\textbf{Protocol} & \textbf{Features} & \textbf{Val} & \textbf{Aro} & \textbf{Leakage} \\
\midrule
Grouped CV (ours)  & DE + HFD       & 66.72\% & 66.25\% & No \\
Pooled random split & DE + BP + Asym & 80.85\% & 81.94\% & Yes \\
Pooled (subset)    & DE + BP + Asym & 82.63\% & 84.29\% & Yes \\
\bottomrule
\end{tabular}
\end{table}

\subsection{Analysis}

The leaky protocol inflates accuracy by \textbf{14--18 percentage points}.
This occurs because:

\begin{enumerate}[noitemsep]
    \item Adjacent 3-second windows from the same 60-second trial are
          autocorrelated---they share nearly identical EEG patterns.
    \item When these correlated windows appear in both training and test sets,
          the model effectively ``memorizes'' trial-specific patterns rather
          than learning generalizable emotional features.
    \item Pooling across subjects further inflates results by mixing
          subject-specific EEG characteristics.
\end{enumerate}

This finding is critical for interpreting published results: \textbf{studies
reporting $>$85\% accuracy on DEAP with classical methods should be scrutinized
for potential data leakage in their evaluation protocol}.


\section{Summary of All Experiments}
\label{sec:exp_summary}

\begin{table}[H]
\centering
\caption{Summary of key experiments and their results.}
\label{tab:all_experiments}
\begin{tabular}{p{4cm}llcc}
\toprule
\textbf{Experiment} & \textbf{Model} & \textbf{Features} & \textbf{Valence} & \textbf{Arousal} \\
\midrule
\multicolumn{5}{l}{\textit{Rigorous evaluation (grouped CV, no leakage)}} \\
\midrule
Paper 13 (grid search) & XGBoost & DE + HFD & \textbf{66.72\%} & \textbf{66.25\%} \\
Paper 13 (default) & XGBoost & DE + HFD & 61.40\% & 62.90\% \\
Focus: DE only & RBF-SVM & DE & 61.88\% & 57.50\% \\
Focus: Phase 2 & HGB & DE & 59.84\% & 56.48\% \\
Ceiling: DE+PLV+MI & RBF-SVM & DE + PLV & 59.06\% & 57.73\% \\
Focus: Soft Vote & Ensemble & DE & 59.38\% & 57.42\% \\
\midrule
\multicolumn{5}{l}{\textit{Leaky evaluation (for reference only)}} \\
\midrule
Leaky XGB & XGBoost & DE+BP+Asym & 80.85\% & 81.94\% \\
Leaky XGB (subset) & XGBoost & DE+BP+Asym & 82.63\% & 84.29\% \\
\bottomrule
\end{tabular}
\end{table}

\subsection{Key Findings}

\begin{enumerate}
    \item \textbf{Best configuration}: XGBoost with DE + HFD features,
          3-second non-overlapping windows, with grid search hyperparameter
          tuning, achieving 66.72\% valence and 66.25\% arousal.

    \item \textbf{Feature importance}: Differential Entropy alone is the most
          effective single feature type. Adding HFD provides complementary
          information. Adding band power, asymmetry, or connectivity features
          does not improve results.

    \item \textbf{Model ranking}: XGBoost $>$ HGB $\approx$ RBF-SVM $>$
          Soft Voting for this task.

    \item \textbf{Window size}: 3-second windows (Paper~13's choice) work well.
          Larger windows reduce sample count without improving accuracy.

    \item \textbf{Feature selection}: Not beneficial for low-dimensional
          feature spaces (e.g., DE with 32 features).

    \item \textbf{Data leakage}: Non-grouped evaluation inflates accuracy by
          15--18 percentage points, explaining the discrepancy with published
          results.

    \item \textbf{Subject variability}: Accuracy varies widely across subjects
          (from $\sim$40\% to $\sim$90\%), reflecting genuine individual
          differences in emotional expressivity through EEG.

    \item \textbf{Hyperparameter tuning}: Grid search improves accuracy by
          $\sim$5 percentage points over default parameters, justifying the
          computational cost of nested CV.
\end{enumerate}


% #############################################################################
% CHAPTER 5: SUMMARY AND CONCLUSION
% #############################################################################
\chapter{Summary and Conclusion}
\label{ch:conclusion}

\section{Summary}

This report presented the design, implementation, and evaluation of a classical
machine learning pipeline for EEG-based emotion recognition using the DEAP
dataset. The key accomplishments are:

\begin{enumerate}
    \item \textbf{Pipeline implementation}: A modular, extensible Python package
          supporting 12 feature types, 10 classifiers, 7 labeling modes, and
          multiple evaluation protocols, with built-in caching and checkpointing.

    \item \textbf{Systematic experimentation}: Over 20 experiments covering
          window size optimization, feature ablation, band selection, feature
          selection, and model comparison, totaling thousands of
          cross-validation folds.

    \item \textbf{Rigorous evaluation}: Consistent use of grouped
          cross-validation (GroupKFold by trial) to prevent data leakage,
          with explicit quantification of leakage effects.

    \item \textbf{Best results}: 66.72\% valence and 66.25\% arousal accuracy
          using XGBoost with Differential Entropy and Higuchi Fractal Dimension
          features under proper evaluation.

    \item \textbf{Practical insights}: Identification of DE as the most
          important feature type, demonstration that feature selection is
          unnecessary for compact feature spaces, and quantification of the
          $\sim$15--18\% accuracy inflation caused by data leakage.
\end{enumerate}


\section{Integration with Real-Time Mobile Application}

The ML pipeline developed in this work is designed to integrate with the
real-time mobile application described in Part~II of this project. The
integration points are:

\begin{itemize}[noitemsep]
    \item \textbf{Feature extraction}: The DE (histogram) and HFD feature
          computation modules can process 3-second EEG windows in real time.
    \item \textbf{Trained model}: The XGBoost classifier with optimized
          hyperparameters can be exported and loaded in the mobile environment.
    \item \textbf{Preprocessing}: The bandpass and notch filter coefficients
          are pre-computed and can be applied as streaming IIR filters.
\end{itemize}

% Placeholder for Part II chapters
% The following chapters will be added for the mobile application component:
% Chapter 6: Mobile Application Architecture
% Chapter 7: Real-Time EEG Acquisition and Processing
% Chapter 8: Mobile Application Implementation
% Chapter 9: System Integration and Testing
% Chapter 10: Overall Conclusions and Future Work


\section{Future Work}

\begin{itemize}
    \item \textbf{Deep learning approaches}: Exploring CNNs, LSTMs, or
          transformer architectures for end-to-end learning from raw EEG.
    \item \textbf{Transfer learning}: Pre-training on large EEG corpora to
          improve cross-subject generalization.
    \item \textbf{Subject adaptation}: Online adaptation techniques that
          calibrate models to individual users with minimal data.
    \item \textbf{Multi-modal fusion}: Combining EEG with peripheral
          physiological signals (GSR, ECG) available in DEAP.
    \item \textbf{Real-time optimization}: Model compression and quantization
          for efficient on-device inference.
\end{itemize}


% ---- REFERENCES ----
\newpage
\begin{thebibliography}{99}
\addcontentsline{toc}{chapter}{References}

\bibitem{koelstra2012deap}
S.~Koelstra, C.~Muhl, M.~Soleymani, J.-S.~Lee, A.~Yazdani, T.~Ebrahimi,
T.~Pun, A.~Nijholt, and I.~Patras,
``DEAP: A Database for Emotion Analysis using Physiological Signals,''
\textit{IEEE Transactions on Affective Computing}, vol.~3, no.~1, pp.~18--31,
Jan. 2012.

\bibitem{russell1980}
J.~A.~Russell,
``A circumplex model of affect,''
\textit{Journal of Personality and Social Psychology}, vol.~39, no.~6,
pp.~1161--1178, 1980.

\bibitem{zheng2015investigating}
W.-L.~Zheng and B.-L.~Lu,
``Investigating critical frequency bands and channels for EEG-based emotion
recognition with deep neural networks,''
\textit{IEEE Transactions on Autonomous Mental Development}, vol.~7, no.~3,
pp.~162--175, 2015.

\bibitem{duan2013differential}
R.-N.~Duan, J.-Y.~Zhu, and B.-L.~Lu,
``Differential entropy feature for EEG-based emotion recognition,''
in \textit{Proc. 6th International IEEE/EMBS Conference on Neural Engineering
(NER)}, pp.~81--84, 2013.

\bibitem{hjorth1970}
B.~Hjorth,
``EEG analysis based on time domain properties,''
\textit{Electroencephalography and Clinical Neurophysiology}, vol.~29, no.~3,
pp.~306--310, 1970.

\bibitem{higuchi1988}
T.~Higuchi,
``Approach to an irregular time series on the basis of the fractal theory,''
\textit{Physica D: Nonlinear Phenomena}, vol.~31, no.~2, pp.~277--283, 1988.

\bibitem{chen2016xgboost}
T.~Chen and C.~Guestrin,
``XGBoost: A scalable tree boosting system,''
in \textit{Proc. 22nd ACM SIGKDD International Conference on Knowledge
Discovery and Data Mining}, pp.~785--794, 2016.

% Add additional references as needed for Papers 11, 13, 14

\end{thebibliography}


% ---- ARABIC SUMMARY ----
% \newpage
% \chapter*{Arabic Summary}
% \addcontentsline{toc}{chapter}{Arabic Summary}
% [Arabic summary to be added here]


\end{document}
